\chapter{Giới thiệu}
\textit{Chương này giới thiệu đề tài nghiên cứu, nêu bật bài toán cần giải quyết, đồng thời trình bày mục tiêu và phạm vi của dự án. Đây là nền tảng để làm rõ động lực và định hướng phát triển một nền tảng học tiếng Anh thông minh.}


\section{Phát biểu bài toán}
Nhu cầu sở hữu các chứng chỉ năng lực tiếng Anh trên toàn cầu đã tăng mạnh, với hàng triệu người học mỗi năm tham gia các kỳ thi chuẩn hóa như IELTS \cite{ielts_official}, TOEIC \cite{toeic_official}, TOEFL và SAT. Tuy nhiên, các phương pháp ôn luyện truyền thống vẫn tồn tại nhiều hạn chế quan trọng:

\begin{itemize}
\item \textbf{Khả năng tiếp cận còn hạn chế}: Các khóa ôn thi chất lượng cao thường đắt đỏ và bị giới hạn về địa lý, khiến nhiều người học khó tiếp cận nguồn tài nguyên phù hợp.
\item \textbf{Thiếu cá nhân hóa}: Tài liệu ôn luyện truyền thống thường áp dụng chung cho mọi đối tượng, chưa phản ánh điểm yếu và phong cách học riêng của từng người học.
\item \textbf{Thiếu cơ hội luyện tập}: Người học thường không có đủ nguồn câu hỏi đa dạng và các bài thi mô phỏng sát thực tế, đặc biệt với kỹ năng Nói.
\item \textbf{Phản hồi chậm}: Các phương pháp đánh giá truyền thống thường trả kết quả chậm, làm giảm tốc độ cải thiện và khả năng học theo vòng lặp.
\item \textbf{Hạn chế trong luyện nói}: Các phần thi Nói như IELTS đòi hỏi sự tương tác trực tiếp, điều rất khó tái hiện trong môi trường tự học. Ngoài ra, việc đánh giá độ trôi chảy khi nói còn cần phân tích khoảng dừng, ngữ điệu và độ chính xác phát âm.
\item \textbf{Khó mở rộng}: Chấm điểm thủ công và hướng dẫn riêng cho từng người không dễ mở rộng để phục vụ số lượng lớn người học.
\item \textbf{Thiếu nhất quán khi triển khai}: Nhiều nền tảng giáo dục gặp tình trạng hoạt động không đồng nhất giữa các môi trường triển khai do lệch cấu hình và xung đột phụ thuộc.
\end{itemize}

Bên cạnh đó, các nền tảng số hiện có còn hạn chế về cách tiếp cận công nghệ, thường thiếu tích hợp AI, chưa bao phủ đầy đủ các dạng bài thi và chưa mang lại trải nghiệm người dùng trực quan. Sự phân mảnh của các công cụ hiện có buộc người học phải sử dụng nhiều nền tảng khác nhau, làm giảm hiệu quả và tăng tải nhận thức. Những vấn đề này tạo ra khoảng trống rõ rệt cho một nền tảng tích hợp, thông minh và dễ tiếp cận, có khả năng hỗ trợ ôn luyện toàn diện cho nhiều định dạng kỳ thi khác nhau đồng thời tận dụng công nghệ AI hiện đại để hỗ trợ việc học \cite{ai_education}.

\section{Mục tiêu}

Mục tiêu chính của dự án là xây dựng một nền tảng học tiếng Anh thông minh giúp người học ôn thi hiệu quả hơn về cả chất lượng lẫn thời gian. Các mục tiêu cụ thể bao gồm:

\begin{enumerate}
    \item \textbf{Hỗ trợ nhiều kỳ thi trên một nền tảng:} Cung cấp một hệ thống thống nhất để người học ôn luyện cho IELTS, TOEIC, TOEFL và SAT. Cách tiếp cận này giảm nhu cầu phải dùng nhiều nguồn khác nhau, đảm bảo định dạng học tập nhất quán và tiết kiệm thời gian.
    
    \item \textbf{Luyện tập tương tác:} Cung cấp các bài luyện Nghe, Nói, Đọc và Viết sát với điều kiện thi thực tế. Người học có thể tăng sự tự tin và độ thành thạo khi luyện tập trong môi trường mô phỏng bài thi thật.
    
    \item \textbf{Phản hồi, chấm điểm tức thời:} Cung cấp phản hồi và chấm điểm kịp thời khi người học hoàn thành bài tập. Hệ thống AI tự động đánh giá bài làm, trả kết quả ngay lập tức và đưa ra nhận xét chi tiết giúp người học nhận biết điểm mạnh, điểm yếu để cải thiện.
    
    \item \textbf{Theo dõi tiến độ học tập:} Cung cấp báo cáo rõ ràng về kết quả và xu hướng tiến bộ theo thời gian. Việc trực quan hóa tiến độ giúp người học xác định khu vực cần cải thiện và phân bổ thời gian học hợp lý hơn.
\end{enumerate}


\section{Phạm vi}

Dự án này giới hạn trong việc thiết kế và phát triển một ứng dụng web có tăng cường AI nhằm hỗ trợ người học chuẩn bị cho các kỳ thi tiếng Anh chuẩn hóa. Phạm vi nghiên cứu xác định ranh giới công việc theo nhóm người dùng mục tiêu, điều kiện vận hành và các ràng buộc của dự án.

\subsection{Đối tượng người dùng}

Người dùng chính của hệ thống là người Việt Nam đang chuẩn bị cho các kỳ thi đánh giá năng lực tiếng Anh như IELTS, TOEFL, TOEIC và SAT, phụ huynh của những học sinh, sinh viên có nhu cầu ôn luyện cho các kỳ thi đánh giá và các giảng viên tiếng Anh.

\subsection{Môi trường vận hành}

Ứng dụng được định hướng vận hành hoàn toàn trực tuyến và có thể truy cập thông qua trình duyệt web trên máy tính để bàn và thiết bị di động. Dự án không giới hạn khu vực địa lý cụ thể, chỉ yêu cầu môi trường sử dụng có kết nối Internet.

\subsection{Thời gian thực hiện}

Phạm vi của dự án bao gồm các hoạt động được thực hiện trong thời gian học phần được giao, gồm phân tích yêu cầu, thiết kế hệ thống, hiện thực, kiểm thử và triển khai. Quá trình bảo trì dài hạn, vận hành thương mại hoặc mở rộng tính năng trong tương lai không nằm trong phạm vi của đề tài.

\subsection{Giới hạn kỹ thuật}

Hệ thống được phát triển theo kiến trúc web nhiều tầng, sử dụng Next.js \cite{nextjs_docs} cho giao diện người dùng, NestJS \cite{nestjs} cho các dịch vụ backend, kết hợp PostgreSQL \cite{postgresql} và Minio \cite{minio} cho lưu trữ dữ liệu và tệp tùy theo miền chức năng. Đối với các tính năng AI, đề tài tập trung vào hai hướng chính: tích hợp các mô hình/API có sẵn cho bài toán hội thoại, phiên âm và tổng hợp tiếng nói; đồng thời xây dựng một mô-đun chấm Writing dựa trên mô hình Phi-3 được fine-tune trên dữ liệu IELTS Writing. Tuy vậy, phạm vi kỹ thuật của đề tài vẫn giới hạn ở mức nguyên mẫu học thuật và triển khai thực nghiệm; các nội dung như huấn luyện mô hình quy mô lớn từ đầu, tối ưu hạ tầng suy luận ở mức sản phẩm thương mại, hay triển khai trên kiến trúc phân tán quy mô lớn không thuộc phạm vi thực hiện.

\subsection{Ràng buộc}

Dự án được thực hiện trong bối cảnh học thuật nên chịu các ràng buộc điển hình như thời gian phát triển hạn chế, tài nguyên có giới hạn và yêu cầu hệ thống phải phù hợp với năng lực triển khai của nhóm bốn thành viên. Vì vậy, việc tăng cường bảo mật ở mức sản phẩm thương mại, cấu hình mức độ sẵn sàng cao và tối ưu hiệu năng không nằm trong mục tiêu chính của đề tài.


\section{Kết luận chương}
Chương này đã trình bày bài toán ôn luyện tiếng Anh chuẩn hóa với những hạn chế của các phương pháp và nền tảng hiện có. Từ đó, nhóm xác định mục tiêu xây dựng một nền tảng tích hợp AI nhằm hỗ trợ nhiều kỳ thi, cung cấp phản hồi tức thì cho kỹ năng Nói và Viết, đồng thời theo dõi tiến độ người học. Phạm vi dự án được giới hạn rõ ràng về đối tượng, thời gian, công nghệ và ràng buộc tài nguyên để đảm bảo tính khả thi trong bối cảnh học thuật.


